\documentclass[11pt,a4paper]{article}

\usepackage[margin=1in]{geometry}

\usepackage{fontspec}
\setmainfont{TeX Gyre Termes}

\usepackage{amsmath}
\usepackage{amssymb}

\usepackage{graphicx}
\usepackage{booktabs}
\usepackage{multirow}
\usepackage{array}

\usepackage{float}
\usepackage{caption}
\usepackage{subcaption}

\usepackage{hyperref}
\usepackage{xcolor}

\usepackage{tikz}
\usetikzlibrary{
    positioning,
    arrows.meta,
    shapes.geometric,
    fit,
    calc
}

\title{
\textbf{FIGNN: Frequency-Informed Graph Representation Learning}\\
for Transferable and Personally Calibrated EEG Workload Estimation
}

\author{
Vignesh Jaisankar\\
Student ID: 24BCE1793
}

\date{September 2026}

\begin{document}

\maketitle

\begin{abstract}
% FINAL ABSTRACT GOES HERE
\end{abstract}

\noindent\textbf{Keywords:}
EEG, graph neural networks, cognitive workload, functional connectivity,
contrastive learning, transfer learning, personal calibration.

\section{Introduction}

\section{Related Work}

\section{Methodology}

\subsection{EEG Graph Representation}

\subsection{Frequency-Informed Graph Neural Network}

\subsection{Contrastive Representation Learning}

\subsection{Stage 1: Leave-One-Subject-Out Evaluation}

\subsection{Stage 2: Subject-Specific Calibration}

\section{Results}

\subsection{Stage 1: Transferability to Unseen Subjects}

\subsection{Stage 2: Personal Calibration}

\section{Discussion}

\section{Limitations}

\section{Conclusion}

\bibliographystyle{plain}
\bibliography{references}

\end{document}
